\textbf{Задача 22} Докажите, что имея источник случайных битов, распределённых равномерно и независимо, можно смоделировать независимые случайные биты, равные единице с вероятностью $p$, если $p$ --- полиномиально вычислимое число, т.~е. по $i$ можно алгоритмически получить $i$-й двоичный бит разложения $p$. Какое среднее число равномерных битов нужно для моделирования одного неравномерного?

\begin{proof}
Пусть $p=0.p_1p_2p_3\ldots$

Возьмем равномерный случай бит $q_1$, если $q_1=0$ и $p_1=1$ то возвращаем $1$,
если $q_1=1$ и $p_1=0$ то возвращаем $0$, иначе берем второй равномерный случай бит $q_2$. Если $q_2=0$ и $p_2=1$ то возвращаем $1$,
если $q_2=1$ и $p_2=0$ то возвращаем $0$, иначе берем третий равномерный случай бит $q_3$, и так далее.
Получили бернуллиевскую случайную величину $\xi$. Найдем её математическое ожидание 
\[\mathbf{E}\,\xi = \sum_{k=1}^{+\infty}p_k\cdot\frac{1}{2^k} = p\]

Найдем математическое ожидания количества случайных битов нужных для генерации $\xi$
\[E = \sum_{k=1}^{+\infty} k \cdot \frac{1}{2^k} = \frac{1}{2}\cdot\left(\frac{1}{1-x}\right)'\left(\frac{1}{2}\right) = \frac{1}{2}\cdot\frac{1}{(1-1/2)^2}=2\]

На самом деле у нас может очень долго не получаться получить значение $\xi$, поэтому можно заканчивать и возвращать $0$ если после $N$ итераций мы еще не вернули ответ, тогда мы будет приближать вероятность $\frac{\left[p\cdot2^N\right]}{2^N}$, что экспоненциально близко к $p$
\end{proof}